% Options for packages loaded elsewhere
% Options for packages loaded elsewhere
\PassOptionsToPackage{unicode}{hyperref}
\PassOptionsToPackage{hyphens}{url}
%
\documentclass[
  ignorenonframetext,
  aspectratio=169,
  russian,
]{beamer}
\newif\ifbibliography
\usepackage{pgfpages}
\setbeamertemplate{caption}[numbered]
\setbeamertemplate{caption label separator}{: }
\setbeamercolor{caption name}{fg=normal text.fg}
\beamertemplatenavigationsymbolshorizontal
% remove section numbering
\setbeamertemplate{part page}{
  \centering
  \begin{beamercolorbox}[sep=16pt,center]{part title}
    \usebeamerfont{part title}\insertpart\par
  \end{beamercolorbox}
}
\setbeamertemplate{section page}{
  \centering
  \begin{beamercolorbox}[sep=12pt,center]{section title}
    \usebeamerfont{section title}\insertsection\par
  \end{beamercolorbox}
}
\setbeamertemplate{subsection page}{
  \centering
  \begin{beamercolorbox}[sep=8pt,center]{subsection title}
    \usebeamerfont{subsection title}\insertsubsection\par
  \end{beamercolorbox}
}
% Prevent slide breaks in the middle of a paragraph
\widowpenalties 1 10000
\raggedbottom
\AtBeginPart{
  \frame{\partpage}
}
\AtBeginSection{
  \ifbibliography
  \else
    \frame{\sectionpage}
  \fi
}
\AtBeginSubsection{
  \frame{\subsectionpage}
}
\usepackage{iftex}
\ifPDFTeX
  \usepackage[T1]{fontenc}
  \usepackage[utf8]{inputenc}
  \usepackage{textcomp} % provide euro and other symbols
\else % if luatex or xetex
  \usepackage{unicode-math} % this also loads fontspec
  \defaultfontfeatures{Scale=MatchLowercase}
  \defaultfontfeatures[\rmfamily]{Ligatures=TeX,Scale=1}
\fi
\usepackage{lmodern}

\ifPDFTeX\else
  % xetex/luatex font selection
\fi
% Use upquote if available, for straight quotes in verbatim environments
\IfFileExists{upquote.sty}{\usepackage{upquote}}{}
\IfFileExists{microtype.sty}{% use microtype if available
  \usepackage[]{microtype}
  \UseMicrotypeSet[protrusion]{basicmath} % disable protrusion for tt fonts
}{}

\usepackage{color}
\usepackage{fancyvrb}
\newcommand{\VerbBar}{|}
\newcommand{\VERB}{\Verb[commandchars=\\\{\}]}
\DefineVerbatimEnvironment{Highlighting}{Verbatim}{commandchars=\\\{\}}
% Add ',fontsize=\small' for more characters per line
\usepackage{framed}
\definecolor{shadecolor}{RGB}{241,243,245}
\newenvironment{Shaded}{\begin{snugshade}}{\end{snugshade}}
\newcommand{\AlertTok}[1]{\textcolor[rgb]{0.68,0.00,0.00}{#1}}
\newcommand{\AnnotationTok}[1]{\textcolor[rgb]{0.37,0.37,0.37}{#1}}
\newcommand{\AttributeTok}[1]{\textcolor[rgb]{0.40,0.45,0.13}{#1}}
\newcommand{\BaseNTok}[1]{\textcolor[rgb]{0.68,0.00,0.00}{#1}}
\newcommand{\BuiltInTok}[1]{\textcolor[rgb]{0.00,0.23,0.31}{#1}}
\newcommand{\CharTok}[1]{\textcolor[rgb]{0.13,0.47,0.30}{#1}}
\newcommand{\CommentTok}[1]{\textcolor[rgb]{0.37,0.37,0.37}{#1}}
\newcommand{\CommentVarTok}[1]{\textcolor[rgb]{0.37,0.37,0.37}{\textit{#1}}}
\newcommand{\ConstantTok}[1]{\textcolor[rgb]{0.56,0.35,0.01}{#1}}
\newcommand{\ControlFlowTok}[1]{\textcolor[rgb]{0.00,0.23,0.31}{\textbf{#1}}}
\newcommand{\DataTypeTok}[1]{\textcolor[rgb]{0.68,0.00,0.00}{#1}}
\newcommand{\DecValTok}[1]{\textcolor[rgb]{0.68,0.00,0.00}{#1}}
\newcommand{\DocumentationTok}[1]{\textcolor[rgb]{0.37,0.37,0.37}{\textit{#1}}}
\newcommand{\ErrorTok}[1]{\textcolor[rgb]{0.68,0.00,0.00}{#1}}
\newcommand{\ExtensionTok}[1]{\textcolor[rgb]{0.00,0.23,0.31}{#1}}
\newcommand{\FloatTok}[1]{\textcolor[rgb]{0.68,0.00,0.00}{#1}}
\newcommand{\FunctionTok}[1]{\textcolor[rgb]{0.28,0.35,0.67}{#1}}
\newcommand{\ImportTok}[1]{\textcolor[rgb]{0.00,0.46,0.62}{#1}}
\newcommand{\InformationTok}[1]{\textcolor[rgb]{0.37,0.37,0.37}{#1}}
\newcommand{\KeywordTok}[1]{\textcolor[rgb]{0.00,0.23,0.31}{\textbf{#1}}}
\newcommand{\NormalTok}[1]{\textcolor[rgb]{0.00,0.23,0.31}{#1}}
\newcommand{\OperatorTok}[1]{\textcolor[rgb]{0.37,0.37,0.37}{#1}}
\newcommand{\OtherTok}[1]{\textcolor[rgb]{0.00,0.23,0.31}{#1}}
\newcommand{\PreprocessorTok}[1]{\textcolor[rgb]{0.68,0.00,0.00}{#1}}
\newcommand{\RegionMarkerTok}[1]{\textcolor[rgb]{0.00,0.23,0.31}{#1}}
\newcommand{\SpecialCharTok}[1]{\textcolor[rgb]{0.37,0.37,0.37}{#1}}
\newcommand{\SpecialStringTok}[1]{\textcolor[rgb]{0.13,0.47,0.30}{#1}}
\newcommand{\StringTok}[1]{\textcolor[rgb]{0.13,0.47,0.30}{#1}}
\newcommand{\VariableTok}[1]{\textcolor[rgb]{0.07,0.07,0.07}{#1}}
\newcommand{\VerbatimStringTok}[1]{\textcolor[rgb]{0.13,0.47,0.30}{#1}}
\newcommand{\WarningTok}[1]{\textcolor[rgb]{0.37,0.37,0.37}{\textit{#1}}}

\usepackage{longtable,booktabs,array}
\newcounter{none} % for unnumbered tables
\usepackage{calc} % for calculating minipage widths
\usepackage{caption}
% Make caption package work with longtable
\makeatletter
\def\fnum@table{\tablename~\thetable}
\makeatother
\usepackage{graphicx}
\makeatletter
\newsavebox\pandoc@box
\newcommand*\pandocbounded[1]{% scales image to fit in text height/width
  \sbox\pandoc@box{#1}%
  \Gscale@div\@tempa{\textheight}{\dimexpr\ht\pandoc@box+\dp\pandoc@box\relax}%
  \Gscale@div\@tempb{\linewidth}{\wd\pandoc@box}%
  \ifdim\@tempb\p@<\@tempa\p@\let\@tempa\@tempb\fi% select the smaller of both
  \ifdim\@tempa\p@<\p@\scalebox{\@tempa}{\usebox\pandoc@box}%
  \else\usebox{\pandoc@box}%
  \fi%
}
% Set default figure placement to htbp
\def\fps@figure{htbp}
\makeatother



\ifLuaTeX
\usepackage[bidi=basic,shorthands=off,provide=*]{babel}
\else
\usepackage[bidi=default,shorthands=off,provide=*]{babel}
\fi
\ifLuaTeX
  \usepackage{selnolig} % disable illegal ligatures
\fi


\setlength{\emergencystretch}{3em} % prevent overfull lines

\providecommand{\tightlist}{%
  \setlength{\itemsep}{0pt}\setlength{\parskip}{0pt}}



 

\usepackage[]{csquotes}

\IfFileExists{plex-otf.sty}{
  %% Full TeXlive
  % \usepackage[%
  %   % math,
  %   RM={Scale=0.94},SS={Scale=0.94},SScon={Scale=0.94},TT={Scale=MatchLowercase,FakeStretch=0.9},DefaultFeatures={Ligatures=Common}
  % ]{plex-otf}
}{
  %% TinyTeX
  \usepackage{libertine}
}

%%% Load theme
% https://deic.uab.cat/~iblanes/beamer_gallery/
\IfFileExists{beamerthemegotham.sty}{
  %% Full TeXlive
  \usetheme{gotham}
  \gothamset{
    numbering=totalpagenumber,
    parttocframe default=off,
    sectiontocframe default=off,
    subsectiontocframe default=off,
  }
}{
  %% TinyTeX
  \usetheme{Madrid}
}
\makeatletter
\@ifpackageloaded{caption}{}{\usepackage{caption}}
\AtBeginDocument{%
\ifdefined\contentsname
  \renewcommand*\contentsname{Содержание}
\else
  \newcommand\contentsname{Содержание}
\fi
\ifdefined\listfigurename
  \renewcommand*\listfigurename{Список иллюстраций}
\else
  \newcommand\listfigurename{Список иллюстраций}
\fi
\ifdefined\listtablename
  \renewcommand*\listtablename{Список таблиц}
\else
  \newcommand\listtablename{Список таблиц}
\fi
\ifdefined\figurename
  \renewcommand*\figurename{Рисунок}
\else
  \newcommand\figurename{Рисунок}
\fi
\ifdefined\tablename
  \renewcommand*\tablename{Таблица}
\else
  \newcommand\tablename{Таблица}
\fi
}
\@ifpackageloaded{float}{}{\usepackage{float}}
\floatstyle{ruled}
\@ifundefined{c@chapter}{\newfloat{codelisting}{h}{lop}}{\newfloat{codelisting}{h}{lop}[chapter]}
\floatname{codelisting}{Список}
\newcommand*\listoflistings{\listof{codelisting}{Листинги}}
\makeatother
\makeatletter
\makeatother
\makeatletter
\@ifpackageloaded{caption}{}{\usepackage{caption}}
\@ifpackageloaded{subcaption}{}{\usepackage{subcaption}}
\makeatother

\usepackage{bookmark}
\IfFileExists{xurl.sty}{\usepackage{xurl}}{} % add URL line breaks if available
\urlstyle{same}
% fallback for those not using the hyperref driver hyperxmp:
\makeatletter
\@ifundefined{xmpquote}{\newcommand{\xmpquote}[1]{#1}}{}
\makeatother
\hypersetup{
  pdftitle={Лабораторная работа №7},
  pdfauthor={Кудякова София Андреевна},
  pdflang={ru},
  hidelinks,
  pdfcreator={LaTeX via pandoc}}


\title{\texorpdfstring{Лабораторная работа №7}{Лабораторная работа №7}}
\subtitle{\texorpdfstring{Дискретно-событийное
моделирование}{Дискретно-событийное моделирование}}
\author{Кудякова София Андреевна}
\date{2026-08-29}

\begin{document}
\frame{\titlepage}

\renewcommand*\contentsname{Содержание}
\begin{frame}[allowframebreaks]
  \frametitle{Содержание}
  \setcounter{tocdepth}{1}
  \tableofcontents
\end{frame}
\setcounter{tocdepth}{1}
\tableofcontents
}

\section{Цель
работы}\label{ux446ux435ux43bux44c-ux440ux430ux431ux43eux442ux44b}

\begin{frame}[fragile]{Цель работы}
\begin{itemize}[<+->]
\tightlist
\item
  Освоить дискретно-событийное моделирование в Julia с использованием
  \texttt{ConcurrentSim}.
\item
  Реализовать модели \textbf{M/M/c} и \textbf{Росса}.
\item
  Провести параметрические эксперименты.
\item
  Построить и проанализировать результаты моделирования.
\item
  Использовать литературное программирование с \texttt{Literate.jl}.
\end{itemize}
\end{frame}

\section{Дискретно-событийное
моделирование}\label{ux434ux438ux441ux43aux440ux435ux442ux43dux43e-ux441ux43eux431ux44bux442ux438ux439ux43dux43eux435-ux43cux43eux434ux435ux43bux438ux440ux43eux432ux430ux43dux438ux435}

\begin{frame}{Дискретно-событийное моделирование}
\textbf{Discrete-Event Simulation (DES)} --- метод моделирования, при
котором состояние системы изменяется в дискретные моменты времени при
наступлении событий.

Основные события:

\begin{itemize}[<+->]
\tightlist
\item
  прибытие заявки;
\item
  начало и окончание обслуживания;
\item
  отказ машины;
\item
  окончание ремонта.
\end{itemize}

Между событиями состояние системы не изменяется.
\end{frame}

\section{Модель M/M/c}\label{ux43cux43eux434ux435ux43bux44c-mmc}

\begin{frame}{Модель M/M/c}
Система \textbf{M/M/c} характеризуется:

\begin{itemize}[<+->]
\tightlist
\item
  пуассоновским входящим потоком с интенсивностью \(\lambda\);
\item
  экспоненциальным временем обслуживания с интенсивностью \(\mu\);
\item
  \(c\) параллельными обслуживающими устройствами.
\end{itemize}

Загрузка системы:

\[
\rho = \frac{\lambda}{c\mu}
\]

Стационарный режим существует при \(\rho < 1\).

Исследовались время ожидания заявок и загрузка очереди.
\end{frame}

\section{Результаты
M/M/c}\label{ux440ux435ux437ux443ux43bux44cux442ux430ux442ux44b-mmc}

\begin{frame}{Результаты M/M/c}
\begin{figure}[H]

{\centering \includegraphics[width=0.85\linewidth,height=\textheight,keepaspectratio]{image/1.png}

}

\caption{Время ожидания заявок}

\end{figure}%

При увеличении нагрузки системы время ожидания заявок возрастает.
\end{frame}

\section{Загрузка очереди
M/M/c}\label{ux437ux430ux433ux440ux443ux437ux43aux430-ux43eux447ux435ux440ux435ux434ux438-mmc}

\begin{frame}{Загрузка очереди M/M/c}
\begin{figure}[H]

{\centering \includegraphics[width=0.85\linewidth,height=\textheight,keepaspectratio]{image/2.png}

}

\caption{Загрузка очереди M/M/c}

\end{figure}%

Рост нагрузки приводит к увеличению вероятности образования очереди.

Параметрическое исследование проводилось при изменении:

\begin{itemize}[<+->]
\tightlist
\item
  количества серверов \(c\);
\item
  интенсивности поступления заявок \(\lambda\);
\item
  загрузки системы \(\rho\).
\end{itemize}

Увеличение числа серверов позволяет снизить нагрузку на очередь.
\end{frame}

\section{Модель
Росса}\label{ux43cux43eux434ux435ux43bux44c-ux440ux43eux441ux441ux430}

\begin{frame}{Модель Росса}
Модель Росса описывает систему с:

\begin{itemize}[<+->]
\tightlist
\item
  \(N\) основными машинами;
\item
  \(S\) резервными машинами;
\item
  ремонтниками;
\item
  отказами и последующим ремонтом.
\end{itemize}

При отказе работающая машина заменяется резервной, если резерв доступен.

Отказавшая машина поступает в ремонт.

При отсутствии резерва система может перейти в состояние краха.
\end{frame}

\section{Динамика исправных
машин}\label{ux434ux438ux43dux430ux43cux438ux43aux430-ux438ux441ux43fux440ux430ux432ux43dux44bux445-ux43cux430ux448ux438ux43d}

\begin{frame}{Динамика исправных машин}
\begin{figure}[H]

{\centering \includegraphics[width=0.85\linewidth,height=\textheight,keepaspectratio]{image/3.png}

}

\caption{Динамика исправных машин}

\end{figure}%

График показывает изменение количества исправных машин во времени
вследствие отказов и ремонтов.
\end{frame}

\section{Распределение времени до
краха}\label{ux440ux430ux441ux43fux440ux435ux434ux435ux43bux435ux43dux438ux435-ux432ux440ux435ux43cux435ux43dux438-ux434ux43e-ux43aux440ux430ux445ux430}

\begin{frame}{Распределение времени до краха}
\begin{figure}[H]

{\centering \includegraphics[width=0.85\linewidth,height=\textheight,keepaspectratio]{image/4.png}

}

\caption{Распределение времени до краха}

\end{figure}%

Гистограмма показывает распределение времени до краха системы по
результатам имитационных экспериментов.

Чем больше время до краха, тем выше живучесть системы.
\end{frame}

\section{Влияние числа
ремонтников}\label{ux432ux43bux438ux44fux43dux438ux435-ux447ux438ux441ux43bux430-ux440ux435ux43cux43eux43dux442ux43dux438ux43aux43eux432}

\begin{frame}{Влияние числа ремонтников}
Исследовалось влияние количества ремонтников на время до краха системы.

Рассматривались два варианта:

\begin{itemize}[<+->]
\tightlist
\item
  \(N=10\) основных машин;
\item
  \(N=6\) основных машин.
\end{itemize}
\end{frame}

\section{Ремонтники: N =
10}\label{ux440ux435ux43cux43eux43dux442ux43dux438ux43aux438-n-10}

\begin{frame}{Ремонтники: N = 10}
\begin{figure}[H]

{\centering \includegraphics[width=0.85\linewidth,height=\textheight,keepaspectratio]{image/5.png}

}

\caption{Влияние числа ремонтников при N=10}

\end{figure}%

Увеличение числа ремонтников позволяет быстрее восстанавливать
отказавшие машины.
\end{frame}

\section{Ремонтники: N =
6}\label{ux440ux435ux43cux43eux43dux442ux43dux438ux43aux438-n-6}

\begin{frame}{Ремонтники: N = 6}
\begin{figure}[H]

{\centering \includegraphics[width=0.85\linewidth,height=\textheight,keepaspectratio]{image/6.png}

}

\caption{Влияние числа ремонтников при N=6}

\end{figure}%

Влияние числа ремонтников зависит от количества основных машин и других
параметров системы.
\end{frame}

\section{Влияние резервных
машин}\label{ux432ux43bux438ux44fux43dux438ux435-ux440ux435ux437ux435ux440ux432ux43dux44bux445-ux43cux430ux448ux438ux43d}

\begin{frame}{Влияние резервных машин}
Количество резервных машин \(S\) является важным параметром модели
Росса.

Чем больше доступный резерв, тем больше отказов можно компенсировать без
немедленного краха системы.

Исследование проведено для \(N=10\) и \(N=6\).
\end{frame}

\section{Резерв: N = 10}\label{ux440ux435ux437ux435ux440ux432-n-10}

\begin{frame}{Резерв: N = 10}
\begin{figure}[H]

{\centering \includegraphics[width=0.85\linewidth,height=\textheight,keepaspectratio]{image/7.png}

}

\caption{Влияние числа резервных машин при N=10}

\end{figure}%

Увеличение количества резервных машин позволяет дольше поддерживать
работоспособность системы после отказов.
\end{frame}

\section{Резерв: N = 6}\label{ux440ux435ux437ux435ux440ux432-n-6}

\begin{frame}{Резерв: N = 6}
\begin{figure}[H]

{\centering \includegraphics[width=0.85\linewidth,height=\textheight,keepaspectratio]{image/8.png}

}

\caption{Влияние числа резервных машин при N=6}

\end{figure}%

График показывает зависимость времени работы системы до краха от
количества резервных машин.
\end{frame}

\section{Параметрические
эксперименты}\label{ux43fux430ux440ux430ux43cux435ux442ux440ux438ux447ux435ux441ux43aux438ux435-ux44dux43aux441ux43fux435ux440ux438ux43cux435ux43dux442ux44b}

\begin{frame}[fragile]{Параметрические эксперименты}
Для автоматизации серии экспериментов использовались:

\begin{Shaded}
\begin{Highlighting}[]
\NormalTok{des\_parameterized.jl}
\NormalTok{ross\_scan\_parameters.jl}
\end{Highlighting}
\end{Shaded}

Результаты сохранялись в CSV-файлы:

\begin{Shaded}
\begin{Highlighting}[]
\NormalTok{data/des\_mmc\_summary.csv}
\NormalTok{data/des\_ross\_summary.csv}
\NormalTok{data/mmc\_summary.csv}
\NormalTok{data/ross\_crash\_times.csv}
\NormalTok{data/ross\_parameter\_scan.csv}
\NormalTok{data/ross\_parameterized.csv}
\NormalTok{data/ross\_summary.csv}
\end{Highlighting}
\end{Shaded}

Параметрический подход позволяет исследовать различные варианты системы
без изменения основной реализации моделей.
\end{frame}

\section{Литературное
программирование}\label{ux43bux438ux442ux435ux440ux430ux442ux443ux440ux43dux43eux435-ux43fux440ux43eux433ux440ux430ux43cux43cux438ux440ux43eux432ux430ux43dux438ux435}

\begin{frame}[fragile]{Литературное программирование}
Для преобразования исходного кода использовался \textbf{Literate.jl}.

В проекте присутствуют:

\begin{Shaded}
\begin{Highlighting}[]
\NormalTok{scripts/des\_literate.jl}
\NormalTok{scripts/des\_parameterized\_literate.jl}
\end{Highlighting}
\end{Shaded}

Литературное программирование объединяет исходный код с поясняющим
текстом и позволяет формировать производные документы.
\end{frame}

\section{Структура
проекта}\label{ux441ux442ux440ux443ux43aux442ux443ux440ux430-ux43fux440ux43eux435ux43aux442ux430}

\begin{frame}[fragile]{Структура проекта}
\begin{Shaded}
\begin{Highlighting}[]
\NormalTok{project/}
\NormalTok{├── data/}
\NormalTok{├── plots/}
\NormalTok{├── scripts/}
\NormalTok{├── src/}
\NormalTok{├── test/}
\NormalTok{├── Project.toml}
\NormalTok{├── Manifest.toml}
\NormalTok{└── README.md}
\end{Highlighting}
\end{Shaded}

\begin{itemize}[<+->]
\tightlist
\item
  \texttt{data/} --- результаты моделирования;
\item
  \texttt{plots/} --- графики;
\item
  \texttt{scripts/} --- скрипты экспериментов.
\end{itemize}
\end{frame}

\section{Выводы}\label{ux432ux44bux432ux43eux434ux44b}

\begin{frame}[fragile]{Выводы}
В ходе лабораторной работы:

\begin{itemize}[<+->]
\tightlist
\item
  изучено дискретно-событийное моделирование;
\item
  реализована модель \textbf{M/M/c};
\item
  исследовано влияние параметров системы на очередь;
\item
  реализована модель \textbf{Росса} с резервированием и ремонтом;
\item
  исследовано влияние числа ремонтников;
\item
  исследовано влияние количества резервных машин;
\item
  выполнены параметрические эксперименты;
\item
  использован \texttt{Literate.jl}.
\end{itemize}

\textbf{Итог:} параметры системы оказывают существенное влияние на время
ожидания, надёжность и время до краха.
\end{frame}




\end{document}
